\begin{theorem}[Integral homology of every toroidal output]\hypertarget{tag-02ZY}{}\label{push:src:A.4.8}\label{push:A48}
\StatementRecord{push:src:A.4.8}{proved}{02ZY}{\SeriesRef{push:src:A.2.4}, \SeriesRef{push:src:A.4.2}, \SeriesRef{push:src:A.4.1}, \SeriesRef{push:src:A.4.7}; \cite[Section 5.6, pp.565--566]{AKMW02AB}; \cite[Theorem~1.6.1, p.~543; Theorem~2.7.1, pp.~551--552]{AKMW02AB}; \cite[Section 2, p.2]{Matsuki00ErrAB}; \cite[Section~5.5, pp.~564--565]{AKMW02AB}}

Let $D_\Sigma$ be the reduced central space of an output of
\ref{push:src:A.2.4} for $u=r^n$. Write $k$ for its number of pure top
ray orbits, $m$ for its number of mixed ray orbits, and
$t_4$ for its number of maximal four-cone orbits. Then
every integral homology group is free, and
\[
 \begin{split}
 (b_0,b_1,b_2,b_3,b_4,b_5,b_6)
   ={}&(1,1,8+k+m,7,\\
       &\hspace{20mm}9k+t_4-2m-1,0,k+m).
 \end{split}
\]
This includes nonreduced central divisors: the homology
is that of their underlying support. The marked nearby
fibre is $X(r^n)$, and its integral monodromy is
$I+nN$ in degree three, where $N$ is the map in \ref{push:src:A.4.2};
it is the identity in every other degree. In particular
the six nonzero Smith factors of this degree-three
monodromy minus the identity are
$n,n,n,n,n,11n$.
\end{theorem}


\begin{proof}
\begin{enumerate}[label=\arabic*.,leftmargin=*]
\item We first compute a reference centre by its normalization
mapping cylinder. Base-change $\mathcal X^+$ by $u=r^n$
and resolve $xy=r^n$ along its smooth conductor $S$.
Put $N_0=\mathcal O_S(m-F)$, using the lattice notation
of \ref{push:src:A.4.1}. The normalization consists of $R^+$ and
$n-1$ copies of $\mathbb P_S(\mathcal O\oplus N_0)$,
in a cycle. The latter assertion follows directly on
the surface-resolution charts: if $x$ has transition
function $g$, the coordinates $x/r^i$ have the same
transition function, while their complementary
coordinates have transition function $g^{-1}$.
Only the closing identification is $T$.

\item In the integral mapping-cylinder sequence the degree-two
map has domain $\mathcal L^n$ and target of rank $11n+1$.
Project each intermediate ruled component to $S$.
The resulting $n-1$ difference blocks are integral
identity matrices. Eliminate them by integral row and
column operations. The remaining map from $\mathcal L$
is $d_2^+$ from \ref{push:src:A.4.1}, together with $n-1$ coordinates
$a\mapsto a(m-F)$. Those coordinates are integral
combinations of the three primitive coordinates of
$d_2^+$. Thus the original degree-two map has primitive
image of rank $10(n-1)+3$, kernel $\mathbb Z^7$, and
cokernel $\mathbb Z^{n+8}$.

\item In degree four, the same elimination uses the $n-1$
identity blocks on the fundamental classes of the bases.
The remaining map from $\mathbb Z$ contains $d_4^+$,
which is primitive by \ref{push:src:A.4.1}; the other entries are the
section differences $m-F$. Hence the degree-four map
is injective with primitive image of rank $n$, and
cokernel of rank $10n+1$. In degree zero it is the
incidence map of an $n$-cycle. In degree six the
normalization contributes one independent fundamental
class per component. Consequently the reference centre
has free integral homology and Betti tuple
\[
 (1,1,n+8,7,10n+1,0,n).
\]

\item On each total
space take the toroidal boundary
\[
 E=\bigl(\{r=0\}\cup f^{-1}(\infty)\bigr)_{\rm red}.
\]
In the cusp charts it is the toric boundary, and outside
the cusp its nontrivial toroidal charts are the plumbing
charts along $S$. The unit transition functions preserve
monomial ideals. \ref{push:src:A.2.4} gives a finite cone complex modulo
the two periods, with the possible self-identifications
of its faces explicitly retained.

\item Make this complex strict by barycentric star subdivisions,
in decreasing cone dimension. At the moment a stratum
is blown up, all of its deeper self-intersections have
already been separated. Its closure is therefore smooth,
and different translates in the chosen orbit are disjoint.
The star at the sum of the primitive basis generators is
the ordinary blow-up of that closure. Repeating over the
finitely many cone orbits gives smooth strict toroidal
embeddings. This is the geometric barycentric subdivision
used in \cite[Section 5.6, pp.565--566]{AKMW02AB}; its
analytic realization by the same monomial ideals is
described in \cite[Section~5.5, pp.~564--565]{AKMW02AB}. Pure $q$-strata
are included in this preparatory step and in the comparison of
the punctured $r$-families.

\item Intersect the two subdivided fans, and take a regular
subdivision, to obtain their explicit common strict
toroidal modification. Apply toroidal weak factorization
to each of its two subdivision maps. The needed theorem
is the finite polyhedral-complex factorization by
nonsingular star subdivisions and inverse subdivisions,
\cite[Theorem~1.6.1, p.~543; Theorem~2.7.1, pp.~551--552]{AKMW02AB}. The corrected
form for a morphism states that its centres are smooth
irreducible closed strata and that all intermediate
complexes map to the target
\cite[Section 2, p.2]{Matsuki00ErrAB}.

\item The complexes are finite, so the cited combinatorial
factorization applies. The analytic
subdivision--modification correspondence of
\cite[Section~5.5, pp.~564--565]{AKMW02AB} realizes each of its finitely
many stars by the indicated monomial blow-up.
The correspondence respects the face identifications
and the bundle units, so the local analytic blow-ups
glue. All intermediate morphisms are proper over the
target of the subdivision map, hence over a sufficiently
small $r$-disc. They remain smooth, and both parameters
remain monomials times units. This gives the required
finite chain of analytic blow-ups and blow-downs
between the total neighbourhoods.

\item We classify the connected centres in this chain. A
nontrivial centre has codimension at least two. For a
cone having a positive $q$-height ray, its period
stabilizer is trivial: a power of a stabilizing element
fixes each of its finitely many primitive rays; the
positive $q$-height first forces the exponent of $G_1$
to vanish and then the exponent of $G_2$ to vanish.
The same argument holds for pure bottom cones. The normalization of the closure of their torus stratum
is the toric variety of its star. If the closed stratum
is smooth, this finite birational normalization map is
an isomorphism.

\item For a mixed cone the centre is contained in $q=r=0$;
its star is complete, and it is a compact smooth toric
point, curve, or surface. It has free integral
cohomology in even degrees only. For a pure top centre,
the cone must have dimension two. Its stabilizer is
precisely $\langle G_1\rangle$, and the chart quotient
and the plumbing identification show that its dense
open part is the corresponding open part of $S$.
The compact smooth centre is therefore a surface
birational to the rational surface $S$. Equivalently,
its proper map to the unmodified top stratum is a
birational map of smooth surfaces. It is a smooth
rational surface and again has free even cohomology.
It need not remain isomorphic to $S$ after earlier
mixed modifications.

\item For a pure bottom centre the cone has dimension two
or three. Its star gives a smooth toric surface or
curve proper over the $r$-disc. The induced character
$r$ is primitive on its quotient lattice, since it
was primitive before quotienting by a smooth cone
on which it vanishes. Its generic fibre is therefore
connected, and is respectively a complete toric
$\mathbb P^1$ or a point. In the surface case the
ordered two-dimensional fan gives a chain of rational
curves as the central support, possibly with positive
multiplicities. In the curve case the support is one
point. Monomial gradient retraction, as in \ref{push:src:A.4.7},
shows that the ordinary integral cohomology of the
centre is free and even, with $H^0=\mathbb Z$ and
$H^4=0$. These centres are not compact; this is why
they must be distinguished from the first two types.

\item For any closed smooth complex centre $C\subset Y$ of
complex codimension $c$, compact or not, the integral
blow-up splitting is
\[
 H^j(\operatorname{Bl}_C Y;\mathbb Z)
 \simeq H^j(Y;\mathbb Z)\oplus
       \bigoplus_{a=1}^{c-1}H^{j-2a}(C;\mathbb Z).
\]
Tubular excision identifies the relative
groups by the integral Thom class; the projective-bundle
theorem for the exceptional divisor supplies the powers
of its hyperplane class. Their exceptional Gysin maps,
together with pullback from $Y$, give the displayed
map. The relation that restriction of the exceptional
class is minus the hyperplane class makes the relative
matrix triangular with unit diagonal. Proper
blow-down pushforward, defined by Borel--Moore duality,
splits the pullback summand. This proves the splitting
also for the noncompact centres just listed.

\item It follows that all odd cohomology groups and all
torsion are unchanged along the factorization. Each
connected centre adds one to $b_2$ and one to the number
of prime boundary components. Each compact centre
adds one to $b_6$ and one to the number of central
prime components: in degree six the only new group
is its top cohomology. A pure bottom centre adds
neither, because its $H^4$ (and, for a curve, its
$H^2$) vanishes. The reverse operations subtract
the same contributions.

\item Consequently $b_2$ minus the number of all boundary
ray orbits, and $b_6$ minus the number of central ray
orbits, are invariants of this comparison. Both
endpoints have exactly one pure bottom ray orbit.
The reference has $n$ central components and
$b_2=n+8$, so the arbitrary output has
\[
 b_2=8+k+m,\qquad b_6=k+m,
 \qquad (b_1,b_3,b_5)=(1,7,0).
\]
All groups are free by the integral splitting in both
directions. The total neighbourhoods retract onto
their central supports by \ref{push:src:A.4.7}, and those compact
analytic supports have finite triangulations. The
universal coefficient theorem therefore transfers
the same freeness and ranks to integral homology.
Finally the independently computed torus-stratum
Euler number $11k+t_4$ of \ref{push:src:A.4.7} gives
$b_4=9k+t_4-2m-1$.

\item Every output agrees with the marked family
$X(r^n)$ off $r=0$. A positive $r$-loop makes $n$
positive $u$-loops. Since $N^2=0$, its monodromy is
$(I+N)^n=I+nN$. Multiplication of the already computed
integral Smith matrix by $n$ gives the stated factors.

\end{enumerate}
\end{proof}
